\chapter{Building It}\label{app:build}

The repository is \file{github.com/adekau/mathbook}. The toolchain is the latest stable Lean 4
(\file{engine/lean-toolchain}, v4.33.1 at this commit); \file{proofs/} pins Mathlib at the
release tag whose toolchain matches, and \file{engine/} depends on \lc{Init} alone, which the
script \file{scripts/check-engine-deps.sh} enforces in CI.

\section*{The engine and its tests}
\begin{minted}{text}
cd engine && lake build && lake test
\end{minted}
\texttt{lake test} replays \file{Tests/golden.tsv}, the wire-level answers recorded by the
differential test against the original TypeScript engine before that engine was deleted, and
the unit tests in \file{Tests/Main.lean}.

\section*{The proofs}
\begin{minted}{text}
cd proofs && lake exe cache get && lake build
\end{minted}
The first command fetches Mathlib's compiled cache; the build then takes a few minutes. Every
top-level theorem's axioms are printed as part of it.

\section*{The notebook}
\begin{minted}{text}
npm install && npm test && npm run bundle
\end{minted}
\texttt{npm test} runs the TypeScript tests, including the native stdio round trip against the
Lean executable. \texttt{npm run bundle} writes the static notebook to
\file{apps/notebook/dist}, with a build stamp on every asset.

\section*{WebAssembly}
\begin{minted}{text}
npm run wasm
\end{minted}
Builds the Lean runtime and \lc{Init} for \texttt{wasm32} from the toolchain's sources (once;
cached under \file{engine/toolchains/}), then links the engine. The first build takes about
ten minutes, most of it elaborating \lc{Init}; the patch the runtime needs is in
\file{engine/wasm/}.

\section*{This book}
\begin{minted}{text}
cd book && latexmk -xelatex -shell-escape show-your-work.tex
\end{minted}
XeLaTeX with \texttt{minted} (Pygments) and \texttt{tcolorbox}. The CI workflow
\file{.github/workflows/build-book.yml} builds it on every push and attaches the PDF to a
release tagged with the date and commit.
